% !TeX program = lualatex
% !TeX encoding = utf8
% !TeX spellcheck = uk_UA
% !TeX root = ../main.tex

%% --------------------------------------------------------
%% Chapter file: Discharge.tex
%% --------------------------------------------------------

\graphicspath{{\currfilebase/Pictures}}

\Opensolutionfile{answer}[\currfilebase/\currfilebase-Answers]
\Writetofile{answer}{\protect\section*{Розділ~\ref*{sec:discharge}. \nameref*{sec:discharge}}}
\chapter{Низькотемпературна плазма та газовий розряд}\label{sec:discharge}\hypertarget{discharge}{}
%% --------------------------------------------------------

Більшість плазм, з~якими людина має справу щодня, зовсім не гарячі: це полум'я дуги
зварювального апарата, світіння люмінесцентної лампи, плазма в~реакторі, що витравлює
транзистори на кремнієвій пластині. Це \emphx{низькотемпературна плазма}\index{Плазма!низькотемпературна}: електрони мають
$T_e\sim1$--$10$~еВ, а~іони та нейтральний газ лишаються майже холодними,
$T_i\approx T_g\sim0{,}03$--$0{,}1$~еВ. Ступінь іонізації~\eqref{eq:ionization-degree} малий,
$10^{-6}$--$10^{-2}$, густина $n_e\sim10^{8}$--$10^{12}~\mathrm{см}^{-3}$. Така плазма
\emph{нерівноважна}: формула Саха~\eqref{eq:saha} до неї не застосовна, бо електрони
нагріває зовнішнє поле, а~енергію вони віддають нейтральному газу повільно (у~пружному
зіткненні електрон передає атому лише частку $\sim2m/M\sim10^{-4}$ своєї енергії).
Плазма живе доти, доки поле підтримує іонізацію, яка компенсує втрати заряджених
частинок на стінки. Цей баланс і~є предметом фізики \emphx{газового розряду}\index{Розряд газовий|textbf}.

У~розділі ми побачимо, як виникає розряд (пробій Таунсенда і~закон Пашена), як
влаштовані тліючий розряд і~дуга, що визначає температуру електронів і~густину плазми
(глобальна модель), і~як працюють сучасні високочастотні джерела плазми для мікроелектроніки.
Детальний виклад є~в~\cite{lieberman2005principles,raizer1991gas,chen2003lecture}.

%% ========================================================
\section{Елементарні процеси: іонізація і~її швидкість}\label{sec:ionization-rate}
%% ========================================================

Основний процес, що підтримує розряд,~--- \emphx{іонізація електронним ударом}\index{Іонізація!електронним ударом|textbf}
$e+A\to A^++2e$. Вона можлива, лише якщо енергія електрона $\varepsilon$ перевищує
потенціал іонізації $\varepsilon_{iz}$ ($13{,}6$~еВ для H, $15{,}8$~еВ для Ar, $24{,}6$~еВ для He).
Класична оцінка Томсона (електрон передає атомному електрону енергію $\Delta\varepsilon$ у~кулонівському
зіткненні, а~іонізація відбувається при $\Delta\varepsilon>\varepsilon_{iz}$) дає переріз
% ~~~~~~~~~~~~~~~~~~~~~~~~~~~~~~~~~~~~~~~~~~~~~~~~~~~~~~~~
\begin{equation}\label{eq:thomson-iz}
	\sigma_{iz}(\varepsilon)=\frac{\pi e^4}{\varepsilon}\left(\frac1{\varepsilon_{iz}}-\frac1\varepsilon\right),
	\qquad \varepsilon>\varepsilon_{iz}.
\end{equation}
% ~~~~~~~~~~~~~~~~~~~~~~~~~~~~~~~~~~~~~~~~~~~~~~~~~~~~~~~~
Переріз лінійно зростає від порога, досягає максимуму
$\sigma_{\max}=\pi e^4/(4\varepsilon_{iz}^2)$ при $\varepsilon=2\varepsilon_{iz}$
(для Ar $\approx3\times10^{-16}~\mathrm{см}^2$, тобто порядку газокінетичного) і~далі спадає як $1/\varepsilon$.
Поряд з~іонізацією відбуваються збудження атомів (саме воно дає світіння розряду),
пружні зіткнення, а~в~електронегативних газах ($\mathrm{O_2}$, $\mathrm{SF_6}$, $\mathrm{CF_4}$)~--- ще й~прилипання
$e+\mathrm{AB}\to\mathrm A+\mathrm B^-$.

Кількість іонізацій в~одиниці об'єму за одиницю часу дорівнює $n_en_gk_{iz}$, де $n_g$~---
густина газу, а~\emphx{константа швидкості}\index{Константа швидкості іонізації|textbf} $k_{iz}=\langle\sigma_{iz}v\rangle$ усереднена за
розподілом електронів. Для максвеллівських електронів з~$T_e\ll\varepsilon_{iz}$ іонізують лише
електрони з~хвоста розподілу. Якщо поблизу порога $\sigma_{iz}\approx C(\varepsilon-\varepsilon_{iz})$, то
(задача~\ref{prb:IonizationRate})
% ~~~~~~~~~~~~~~~~~~~~~~~~~~~~~~~~~~~~~~~~~~~~~~~~~~~~~~~~
\begin{equation}\label{eq:kiz}
	\tcbhighmath{k_{iz}(T_e)=C\bar v_e\,(\varepsilon_{iz}+2T_e)\,e^{-\varepsilon_{iz}/T_e},\qquad \bar v_e=\sqrt{\frac{8T_e}{\pi m}}.}
\end{equation}
% ~~~~~~~~~~~~~~~~~~~~~~~~~~~~~~~~~~~~~~~~~~~~~~~~~~~~~~~~
Головне в~\eqref{eq:kiz}~--- арреніусівський множник $e^{-\varepsilon_{iz}/T_e}$. Для аргону
добре працює апроксимація $k_{iz}\approx2{,}3\times10^{-8}\,T_e^{0{,}59}e^{-17{,}4/T_e}~\mathrm{см}^3/\mathrm{с}$
($T_e$ в~еВ)~\cite{lieberman2005principles}: при $T_e=2$~еВ $k_{iz}\approx2\times10^{-11}$, а~при
$T_e=4$~еВ уже $\approx4\times10^{-10}~\mathrm{см}^3/\mathrm{с}$. Подвоєння температури збільшує
швидкість іонізації в~двадцять разів. Через таку різку залежність $T_e$ у~розряді ледь змінюється
(див. розд.~\ref{sec:global-model}): невелика зміна $T_e$ компенсує зміну втрат у~кілька разів.

\emph{Оцінка.} Втрати енергії електрона в~пружних зіткненнях малі ($2m/M$ за зіткнення), тому
основний \enquote{стік} енергії електронів~--- непружні процеси. На одну іонізацію в~аргоні
при $T_e\sim3$~еВ припадає кілька збуджень, і~повна енергія, яку витрачає електронний газ на
створення однієї пари електрон--іон, становить $\varepsilon_c\approx50$--$60$~еВ, у~кілька разів більше
за $\varepsilon_{iz}$. Більша частина цієї енергії виходить з~плазми у~вигляді випромінювання.

%% ========================================================
\section{Пробій газу: механізм Таунсенда і~закон Пашена}\index{Таунсенда механізм|textbf}\index{Пробій газу|textbf}\label{sec:townsend}
%% ========================================================

Нехай між плоскими електродами на відстані $d$ прикладено напругу $V$, поле $E=V/d$.
Випадковий вільний електрон (від космічного випромінювання чи фотоефекту) прискорюється
до анода і~іонізує газ. Число іонізацій на одиницю довжини шляху~--- \emphx{перший
коефіцієнт Таунсенда} $\alpha$. Оцінимо його. Між зіткненнями електрон проходить у~середньому
$\lambda=1/(n_g\sigma)\propto1/p$ і~набирає енергію $eE\lambda$. Імовірність пройти без
зіткнення відстань $\varepsilon_{iz}/(eE)$, потрібну для набору енергії іонізації, дорівнює
$\exp[-\varepsilon_{iz}/(eE\lambda)]$; зіткнень на одиницю довжини $1/\lambda$. Отже,
% ~~~~~~~~~~~~~~~~~~~~~~~~~~~~~~~~~~~~~~~~~~~~~~~~~~~~~~~~
\begin{equation}\label{eq:townsend-alpha}
	\alpha=Ap\,\exp\left(-\frac{Bp}{E}\right),\qquad A=\frac1{\lambda p},\quad B=\frac{\varepsilon_{iz}A}{e},
\end{equation}
% ~~~~~~~~~~~~~~~~~~~~~~~~~~~~~~~~~~~~~~~~~~~~~~~~~~~~~~~~
де $A$ і~$B$~--- сталі газу (на практиці їх беруть з~експерименту: для повітря
$A\approx15~\mathrm{см}^{-1}\mathrm{Торр}^{-1}$, $B\approx365$~В/(см$\,\cdot\,$Торр); для аргону $A\approx12$,
$B\approx180$). Важливо, що $\alpha/p$ залежить лише від $E/p$~--- енергії, набраної на довжині пробігу.

Один електрон, що вилетів з~катода, дає на аноді лавину $e^{\alpha d}$ електронів; утворені
$e^{\alpha d}-1$ іонів повертаються на катод і~вибивають з~нього $\gamma(e^{\alpha d}-1)$ нових
електронів, де $\gamma\sim10^{-2}$--$10^{-1}$~--- \emphx{коефіцієнт вторинної емісії}\index{Вторинної емісії коефіцієнт|textbf} (другий
коефіцієнт Таунсенда). Якщо кожен вихідний електрон відтворює хоча б одного наступника,
розряд стає \emph{самостійним} і~не потребує зовнішнього іонізатора:
% ~~~~~~~~~~~~~~~~~~~~~~~~~~~~~~~~~~~~~~~~~~~~~~~~~~~~~~~~
\begin{equation}\label{eq:townsend-criterion}
	\gamma\left(e^{\alpha d}-1\right)=1\quad\Longleftrightarrow\quad \alpha d=\ln\left(1+\frac1\gamma\right).
\end{equation}
% ~~~~~~~~~~~~~~~~~~~~~~~~~~~~~~~~~~~~~~~~~~~~~~~~~~~~~~~~
Підставивши~\eqref{eq:townsend-alpha} з~$E=V/d$, отримуємо напругу пробою~--- \emphx{закон Пашена}\index{Пашена закон|textbf}:
% ~~~~~~~~~~~~~~~~~~~~~~~~~~~~~~~~~~~~~~~~~~~~~~~~~~~~~~~~
\begin{equation}\label{eq:paschen}
	\tcbhighmath{V_b=\frac{B\,pd}{\ln(A\,pd)-\ln\ln(1+1/\gamma)}.}
\end{equation}
% ~~~~~~~~~~~~~~~~~~~~~~~~~~~~~~~~~~~~~~~~~~~~~~~~~~~~~~~~
Напруга пробою залежить не від $p$ і~$d$ окремо, а~лише від добутку $pd$ (закон подібності).
Крива Пашена (рис.~\ref{fig:paschen}) має мінімум при
$(pd)_{\min}=e\ln(1+1/\gamma)/A$, де $V_{\min}=B\,(pd)_{\min}$ (задача~\ref{prb:Paschen}); для повітря
$V_{\min}\approx300$~В при $pd\approx0{,}8$~Торр$\,\cdot\,$см.
Праворуч від мінімуму електрон зазнає багато зіткнень, але за одне зіткнення набирає замало енергії,
тому потрібне сильне поле. Ліворуч від мінімуму зіткнень замало, щоб розвинулася лавина: електрон
долітає до анода, майже нікого не іонізувавши. Звідси практичні наслідки:
повітряний проміжок $1$~см при атмосферному тиску ($pd=760$~Торр$\,\cdot\,$см) пробивається при
$\approx30$~кВ; вакуумні вимикачі і~високовольтні вводи працюють на лівій гілці, а~пробій усередині
тонкої трубки при низькому тиску може піти \enquote{довгим шляхом} уздовж трубки, а~не коротким
поперек.

\begin{figure}[ht]
	\centering
	\localinput{paschen.tikz}
	\caption{Криві Пашена~\eqref{eq:paschen} для повітря ($A=15$, $B=365$) і~аргону ($A=12$, $B=180$) при
	$\gamma=0{,}01$. Точками позначено мінімуми $V_{\min}=eB\ln(1+1/\gamma)/A$.}
	\label{fig:paschen}
\end{figure}

Механізм Таунсенда описує пробій, поки лавина не спотворює зовнішнього поля. При
$pd\gtrsim200$~Торр$\,\cdot\,$см (довгі іскри, блискавка) просторовий заряд головки лавини
сам створює поле, порівнянне з~прикладеним. Тоді розряд поширюється як \emphx{стример}\index{Стример|textbf}~---
тонкий іонізаційний фронт, що рухається зі швидкістю $10^{7}$--$10^{8}$~см/с і~підживлюється
фотоіонізацією перед головкою, а~катод і~коефіцієнт $\gamma$ ролі вже не грають.

%% ========================================================
\section{Тліючий розряд і~дуга}\index{Розряд!тліючий|textbf}\label{sec:glow}
%% ========================================================

Після пробою при обмеженому струмі (зовнішній резистор) встановлюється \emphx{тліючий
розряд}. Його найпомітніша риса~--- шарувата будова (рис.~\ref{fig:glow}). Біля катода
лежить \emphx{катодне падіння}\index{Катодне падіння|textbf}: шар завтовшки в~кілька довжин пробігу, на якому спадає майже
вся напруга, $V_c\sim100$--$400$~В (порядку $V_{\min}$ кривої Пашена). У~ньому іони прискорюються
до катода, вибивають вторинні електрони, і~ті розмножуються лавиною, як у~механізмі Таунсенда.
Прискорені електрони виходять із~шару в~\emph{негативне світіння}~--- область майже без поля,
де вони втрачають енергію на збудження та іонізацію. Далі, у~\emph{фарадеєвому темному
просторі}, швидких електронів уже немає, а~поле ще мале. Решту проміжку до анода заповнює
\emphx{позитивний стовп}\index{Позитивний стовп|textbf}~--- однорідна квазінейтральна плазма зі слабким поздовжнім полем
$E\sim1$--$10$~В/см.

\begin{figure}[ht]
	\centering
	\localinput{glow-structure.tikz}
	\caption{Будова тліючого розряду в~трубці (зверху) і~розподіл потенціалу $\varphi(x)$ (знизу):
	1~--- астонів темний простір, 2~--- катодне світіння, 3~--- катодний (круксів) темний простір,
	4~--- негативне світіння, 5~--- фарадеїв темний простір, 6~--- позитивний стовп, 7~--- анодне
	світіння. Майже вся напруга спадає в~катодному шарі (1--3).}
	\label{fig:glow}
\end{figure}

\subsection*{Позитивний стовп: температура з~балансу частинок}

Позитивний стовп~--- найпростіший приклад самоузгодженої плазми. Заряджені частинки
дифундують до стінок трубки радіуса $R$ амбіполярно (розд.~\ref{sec:ambipolar}) і~там
рекомбінують; втрати компенсує іонізація в~об'ємі. Рівняння неперервності для густини,
$D_a\nabla^2n+n_gk_{iz}n=0$, з~умовою $n(R)\approx0$ має розв'язок $n(r)=n_0J_0(2{,}405\,r/R)$ і~дає
\emph{умову Шотткі}
% ~~~~~~~~~~~~~~~~~~~~~~~~~~~~~~~~~~~~~~~~~~~~~~~~~~~~~~~~
\begin{equation}\label{eq:schottky}
	n_gk_{iz}(T_e)=D_a\left(\frac{2{,}405}{R}\right)^2,\qquad D_a\approx\frac{T_e}{M\nu_{in}}.
\end{equation}
% ~~~~~~~~~~~~~~~~~~~~~~~~~~~~~~~~~~~~~~~~~~~~~~~~~~~~~~~~
Тут $D_a\approx D_i(1+T_e/T_i)\approx T_e/(M\nu_{in})$ з~\eqref{eq:ambipolar}, а~$\nu_{in}\propto n_g$~---
частота зіткнень іонів з~атомами. Оскільки $n_g\propto p$, умова~\eqref{eq:schottky} пов'язує
$T_e$ лише з~параметром подібності $pR$: \emph{густина плазми з~неї випадає зовсім}. Температуру
електронів задає геометрія і~тиск, а~струм розряду (через баланс енергії) визначає лише густину.
Зі зменшенням $pR$ дифузійні втрати зростають, і~$T_e$ зростає, щоб іонізація встигала їх
компенсувати. Цю ж~логіку ми повторимо в~глобальній моделі (розд.~\ref{sec:global-model}).

\subsection*{Вольт-амперна характеристика і~дуга}

Зі зростанням струму розряд проходить кілька режимів. При малих струмах ($\lesssim10^{-6}$~А)
горить \emph{темний таунсендівський} розряд. Далі напруга стрибком падає до значення
\emph{нормального} тліючого розряду. У~ньому катодне падіння і~густина струму на катоді
сталі ($j/p^2=\mathrm{const}$), а~зростання струму лише збільшує вкриту світінням площу катода.
Коли світіння вкриває весь катод, настає \emph{аномальний} тліючий розряд: щоб збільшити струм,
треба підвищувати $V_c$.
При струмах $\gtrsim1$~А іонне бомбардування розігріває катод до термоемісії (або на ньому
утворюються мікроскопічні катодні плями з~автоемісією), і~розряд переходить у~\emphx{дугу}\index{Дуга|textbf}.
Катодне падіння в~дузі лише $\sim10$~В, бо електрони емітує сам гарячий катод, а~не
вибивають іони. Напруга на дузі падає зі зростанням струму, тобто дуга має \emph{спадну}
характеристику, тому без баластного опору вона нестійка. Дуга атмосферного тиску (зварювання,
плазмотрони, вимикачі) близька до локальної термодинамічної рівноваги: $T_e\approx T_i\approx T_g\sim1$~еВ,
бо часті зіткнення вирівнюють температури, і~до неї вже можна застосувати формулу Саха.

%% ========================================================
\section{Приелектродний шар у~розряді}\index{Шар біля стінки}\label{sec:discharge-sheath}
%% ========================================================

Будь-яка розрядна плазма відділена від стінок і~електродів шаром позитивного
просторового заряду. Його структуру для малого перепаду потенціалу ми вже розглянули
(розд.~\ref{sec:bohm-sheath}): іони входять у~шар зі швидкістю Бома $u_B=\sqrt{T_e/M}$~\eqref{eq:bohm-criterion},
а~густина на межі шару внаслідок прискорення в~передшарі становить $n_s\approx e^{-1/2}n_0\approx0{,}6\,n_0$.
Отже, потік іонів на поверхню
% ~~~~~~~~~~~~~~~~~~~~~~~~~~~~~~~~~~~~~~~~~~~~~~~~~~~~~~~~
\begin{equation}\label{eq:bohm-flux}
	\Gamma_i\approx0{,}6\,n_0u_B
\end{equation}
% ~~~~~~~~~~~~~~~~~~~~~~~~~~~~~~~~~~~~~~~~~~~~~~~~~~~~~~~~
\emph{не залежить від потенціалу поверхні}: його задає плазма, а~не стінка. Ця обставина
лежить в~основі всієї технології плазмової обробки.

На ізольовану стінку іони падають з~енергією $e|\Delta\varphi_f|+T_e/2$, де
$\Delta\varphi_f$~--- плаваючий потенціал~\eqref{eq:floating-potential}; для аргону це $\approx5{,}2\,T_e$,
тобто $\sim15$~еВ. Щоб отримати іони з~енергією в~сотні електронвольт, на електрод подають
великий від'ємний потенціал $-V_0$, $eV_0\gg T_e$. Тоді електрони в~шарі практично відсутні, і~шар
містить лише іони, що прискорюються. Нехай $\Phi(x)>0$~--- спад потенціалу від межі шару ($x=0$)
вглиб, $J=en_iv_i$~--- стала густина іонного струму, $v_i=\sqrt{2e\Phi/M}$ (початковою швидкістю
$u_B$ нехтуємо). Рівняння Пуассона
% ~~~~~~~~~~~~~~~~~~~~~~~~~~~~~~~~~~~~~~~~~~~~~~~~~~~~~~~~
\begin{equation*}
	\frac{d^2\Phi}{dx^2}=4\pi en_i=\frac{4\pi J}{\sqrt{2e\Phi/M}}
\end{equation*}
% ~~~~~~~~~~~~~~~~~~~~~~~~~~~~~~~~~~~~~~~~~~~~~~~~~~~~~~~~
після множення на $d\Phi/dx$ і~інтегрування з~$\Phi=\Phi'=0$ при $x=0$ дає $\Phi^{3/4}\propto x$,
а~умова $\Phi(s)=V_0$ на електроді~--- \emphx{закон Чайлда--Ленгмюра}\index{Чайлда--Ленгмюра закон|textbf}:
% ~~~~~~~~~~~~~~~~~~~~~~~~~~~~~~~~~~~~~~~~~~~~~~~~~~~~~~~~
\begin{equation}\label{eq:child-langmuir}
	\tcbhighmath{J=\frac1{9\pi}\sqrt{\frac{2e}{M}}\,\frac{V_0^{3/2}}{s^2}.}
\end{equation}
% ~~~~~~~~~~~~~~~~~~~~~~~~~~~~~~~~~~~~~~~~~~~~~~~~~~~~~~~~
Той самий закон \enquote{трьох других} обмежує струм електронів у~вакуумному діоді
(з~заміною $M\to m$). У~плазмі струм $J$ задано наперед потоком Бома, $J=e\,n_su_B$, тож
\eqref{eq:child-langmuir} визначає не струм, а~\emph{товщину шару}:
% ~~~~~~~~~~~~~~~~~~~~~~~~~~~~~~~~~~~~~~~~~~~~~~~~~~~~~~~~
\begin{equation}\label{eq:sheath-thickness}
	s=\frac{\sqrt2}{3}\,r_D\left(\frac{2eV_0}{T_e}\right)^{3/4},
\end{equation}
% ~~~~~~~~~~~~~~~~~~~~~~~~~~~~~~~~~~~~~~~~~~~~~~~~~~~~~~~~
де $r_D$~--- дебаївський радіус~\eqref{eq:debye-radius} при густині $n_s$. Шар під високою напругою
в~десятки разів товщий за дебаївський: при $T_e=3$~еВ, $n_s=10^{10}~\mathrm{см}^{-3}$ і~$V_0=300$~В
маємо $r_D\approx0{,}13$~мм і~$s\approx3$~мм (задача~\ref{prb:ChildLangmuir}).

\emph{Енергія іонів на поверхні.} Якщо шар беззіткненнєвий ($\lambda_i\gg s$), кожен іон приходить
на електрод з~енергією $eV_0$, а~розподіл за енергіями вузький. Кут розходження
визначає відношення поперечної теплової швидкості до набраної в~шарі:
$\theta\approx\sqrt{T_i/(2eV_0)}\sim1^\circ$. Саме цей майже нормальний пучок іонів робить
можливим анізотропне травлення (розд.~\ref{sec:plasma-processing}). Якщо ж~$\lambda_i<s$ (тиск
$\gtrsim100$~мТорр), іони перезаряджаються в~шарі, і~їхній енергетичний спектр розмивається в~бік
малих енергій, а~кутовий розширюється.

%% ========================================================
\section{Глобальна модель розряду}\label{sec:global-model}
%% ========================================================

Як для заданого джерела (газ, тиск, розміри, потужність) оцінити $T_e$ і~$n_0$, не
розв'язуючи рівнянь переносу? \emphx{Глобальна модель}\index{Глобальна модель розряду|textbf} усереднює плазму за об'ємом $V$ і~
записує два балансові співвідношення.

\emph{Баланс частинок.} Іонізація в~об'ємі компенсує потік Бома на поверхню площею $A$:
% ~~~~~~~~~~~~~~~~~~~~~~~~~~~~~~~~~~~~~~~~~~~~~~~~~~~~~~~~
\begin{equation*}
	n_gk_{iz}n_0V=n_0u_BA_{\rm eff},
\end{equation*}
% ~~~~~~~~~~~~~~~~~~~~~~~~~~~~~~~~~~~~~~~~~~~~~~~~~~~~~~~~
де ефективна площа $A_{\rm eff}=h A$ враховує, що густина біля стінки менша, ніж у~центрі
(при низькому тиску $h\sim0{,}3$--$0{,}5$, з~ростом тиску $h$ спадає як $\lambda_i^{1/2}$). Густина $n_0$ скорочується, і~ми
отримуємо рівняння лише для $T_e$:
% ~~~~~~~~~~~~~~~~~~~~~~~~~~~~~~~~~~~~~~~~~~~~~~~~~~~~~~~~
\begin{equation}\label{eq:global-te}
	\tcbhighmath{\frac{k_{iz}(T_e)}{u_B(T_e)}=\frac1{n_gd_{\rm eff}},\qquad d_{\rm eff}=\frac V{A_{\rm eff}}.}
\end{equation}
% ~~~~~~~~~~~~~~~~~~~~~~~~~~~~~~~~~~~~~~~~~~~~~~~~~~~~~~~~
Це той самий висновок, що й~для позитивного стовпа: \emph{температуру електронів задає баланс
частинок, вона залежить тільки від $n_gd_{\rm eff}$ і~не залежить від потужності}. Через
арреніусівську залежність $k_{iz}(T_e)$ розв'язок лежить у~вузькій смузі $T_e\approx2$--$5$~еВ
майже для всіх джерел: малі розряди при низькому тиску мають вищу $T_e$, великі при високому~--- нижчу.

\emph{Баланс енергії.} Кожна пара електрон--іон, що покидає плазму, забирає енергію
% ~~~~~~~~~~~~~~~~~~~~~~~~~~~~~~~~~~~~~~~~~~~~~~~~~~~~~~~~
\begin{equation}\label{eq:energy-per-pair}
	\mathcal E_T=\varepsilon_c+2T_e+\varepsilon_i,
\end{equation}
% ~~~~~~~~~~~~~~~~~~~~~~~~~~~~~~~~~~~~~~~~~~~~~~~~~~~~~~~~
де $\varepsilon_c$~--- енергія, витрачена в~зіткненнях на створення пари (розд.~\ref{sec:ionization-rate}),
$2T_e$~--- середня кінетична енергія електрона, що долає потенціальний бар'єр шару
(потік через бар'єр збагачений швидкими частинками), а~$\varepsilon_i\approx e\Delta\varphi+T_e/2$~--- енергія, яку
іон набирає в~передшарі й~шарі. Поглинена плазмою потужність $P_{\rm abs}$ виходить з~неї разом
з~потоком частинок: $P_{\rm abs}=e\,n_0u_BA_{\rm eff}\,\mathcal E_T$ (тут $\mathcal E_T$ у~вольтах). Звідси
% ~~~~~~~~~~~~~~~~~~~~~~~~~~~~~~~~~~~~~~~~~~~~~~~~~~~~~~~~
\begin{equation}\label{eq:global-n}
	\tcbhighmath{n_0=\frac{P_{\rm abs}}{e\,u_BA_{\rm eff}\,\mathcal E_T}.}
\end{equation}
% ~~~~~~~~~~~~~~~~~~~~~~~~~~~~~~~~~~~~~~~~~~~~~~~~~~~~~~~~
Густина пропорційна потужності, а~$T_e$ від неї не залежить~--- розподіл ролей, підтверджений
вимірюваннями в~багатьох джерелах.

\begin{Estimate}
Аргон, $p=10$~мТорр ($n_g\approx3{,}2\times10^{14}~\mathrm{см}^{-3}$ при кімнатній
температурі), $d_{\rm eff}=10$~см. Тоді $n_gd_{\rm eff}\approx3{,}2\times10^{15}~\mathrm{см}^{-2}$,
і~розв'язок~\eqref{eq:global-te} дає $T_e\approx2{,}8$~еВ ($u_B\approx2{,}6\times10^5$~см/с). При
$\varepsilon_c\approx60$~еВ і~$\varepsilon_i\approx5{,}2\,T_e$ маємо $\mathcal E_T\approx80$~еВ, і~потужність
$500$~Вт на $A_{\rm eff}=10^3~\mathrm{см}^2$ підтримує $n_0\approx1{,}5\times10^{11}~\mathrm{см}^{-3}$
(задача~\ref{prb:GlobalModel}). Це типові параметри індуктивного джерела.
\end{Estimate}

%% ========================================================
\section{Високочастотні джерела плазми}\label{sec:rf-sources}
%% ========================================================

Розряд постійного струму має два недоліки: електроди контактують з~плазмою (і~розпорошуються),
а~на діелектричних поверхнях (пластина з~оксидом) накопичується заряд. Тому в~технології
використовують високочастотні (ВЧ) розряди. Стандартна частота $13{,}56$~МГц лежить між
іонною і~електронною плазмовими частотами: $\omega_{pi}<\omega\ll\omega_{pe}$. Електрони
миттєво відстежують поле, а~іони реагують лише на усереднений за періодом потенціал шару. Основні
типи джерел показано на рис.~\ref{fig:rf-sources}, їхні параметри порівняно в~табл.~\ref{tab:rf-sources}.

\begin{figure}[ht]
	\centering
	\localinput{rf-sources.tikz}
	\caption{Схеми ВЧ-джерел плазми: (а)~ємнісне (ЄЗП), енергію передають коливні шари біля
	електродів; (б)~індуктивне (ІЗП), вихрове поле котушки проникає в~плазму на глибину скін-шару
	$\delta$; (в)~ЕЦР/геліконне, мікрохвилі або геліконна хвиля поглинаються в~магнітному полі.
	В~(б) окреме ВЧ-зміщення підкладки керує енергією іонів незалежно від густини.}
	\label{fig:rf-sources}
\end{figure}

\subsection*{Ємнісний розряд}

У~\emphx{ємнісно-зв'язаній плазмі}\index{Розряд!ємнісний|textbf} (ЄЗП) ВЧ-напругу подають на один з~двох плоских
електродів через розділювальний конденсатор. Плазма має потенціал, вищий за потенціал
обох електродів, тож біля кожного утворюється шар, товщина якого коливається з~частотою поля.
Електрони нагріваються омічно (зіткнення в~об'ємі) і~\emph{стохастично}: вони відбиваються від
рухомої межі шару, як м'ячик від ракетки, і~набирають енергію. Оскільки електрони значно
рухливіші за іони, сумарний струм на електрод за період може дорівнювати нулю лише за умови, що
конденсатор зарядиться і~на електроді виникне від'ємне \emphx{автозміщення}\index{Автозміщення|textbf}. Якщо площі
електродів різні, більший спад напруги припадає на шар біля меншого електрода:
$V_1/V_2\approx(A_2/A_1)^q$, ідеалізована модель дає $q=4$, а~на практиці $q\approx1$--$2$. Тому
пластину кладуть на менший (ВЧ) електрод, а~більшим є заземлені стінки камери. Головна вада ЄЗП:
однією напругою визначаються і~густина плазми ($n\sim10^{9}$--$10^{10}~\mathrm{см}^{-3}$), і~енергія іонів
(сотні еВ), тож незалежно керувати ними неможливо. Частковий вихід дають двочастотні ЄЗП
(десятки МГц для густини, $\sim2$~МГц для енергії іонів).

\subsection*{Індуктивний розряд}

У~\emphx{індуктивно-зв'язаній плазмі}\index{Розряд!індуктивний|textbf} (ІЗП) ВЧ-струм тече в~котушці над діелектричним
вікном. Змінне магнітне поле котушки наводить у~плазмі вихрове електричне поле $E_\theta$, яке
нагріває електрони, як вторинну обмотку трансформатора. Поле проникає в~плазму лише на
глибину скін-шару: у~беззіткненнєвому випадку з~провідності~\eqref{eq:drude} при $\nu\ll\omega$
$\delta=c/\omega_{pe}$, для $n=10^{11}~\mathrm{см}^{-3}$ $\delta\approx1{,}7$~см. Тут немає
високовольтних шарів біля котушки, тож густина вища ($10^{11}$--$10^{12}~\mathrm{см}^{-3}$ при
$1$--$50$~мТорр), а~енергією іонів керує окреме ВЧ-зміщення столика з~пластиною. При малій
потужності ІЗП працює в~ємнісному режимі (E-режим) через напругу на котушці і~зі зростанням
потужності стрибком переходить в~індуктивний (H-режим) з~густиною на порядок вищою.

\subsection*{ЕЦР- і~геліконні джерела}

У~магнітному полі $B$ мікрохвилі поглинаються резонансно, де $\omega=\omega_{ce}$
(розд.~\ref{sec:cold-magnetized}). Для стандартної частоти $2{,}45$~ГГц резонансне поле
$B=875$~Гс. Звичайна хвиля не проходить у~плазму з~густиною, вищою за критичну
$n_c=m\omega^2/(4\pi e^2)\approx7{,}4\times10^{10}~\mathrm{см}^{-3}$, а~права поляризована хвиля, введена
з~боку сильного поля ($\omega<\omega_{ce}$), поширюється як свистова хвиля~\eqref{eq:whistler} і~не
має обмеження за густиною, доки не дійде до резонансу. Так \emphx{ЕЦР-джерело}\index{Джерело плазми!ЕЦР|textbf} дає
$n\sim10^{11}$--$10^{12}~\mathrm{см}^{-3}$ при тиску лише $0{,}1$--$1$~мТорр.

\emphx{Геліконне джерело}\index{Джерело плазми!геліконне|textbf} використовує ту саму свистову гілку на ВЧ-частотах
($\omega_{ci}\ll\omega\ll\omega_{ce}$) в~обмеженому циліндрі. При $\omega_{pe}^2\gg\omega\omega_{ce}$
закон дисперсії свистової хвилі, що поширюється під кутом $\theta$ до поля,
$c^2k^2/\omega^2\approx\omega_{pe}^2/(\omega\omega_{ce}\cos\theta)$, з~$k\cos\theta=k_z$ набуває вигляду
% ~~~~~~~~~~~~~~~~~~~~~~~~~~~~~~~~~~~~~~~~~~~~~~~~~~~~~~~~
\begin{equation}\label{eq:helicon}
	kk_z=\frac{\omega\,\omega_{pe}^2}{\omega_{ce}c^2}=\frac{4\pi ne\,\omega}{Bc}
	\quad\Longrightarrow\quad n=\frac{Bkk_zc}{4\pi e\omega}.
\end{equation}
% ~~~~~~~~~~~~~~~~~~~~~~~~~~~~~~~~~~~~~~~~~~~~~~~~~~~~~~~~
При фіксованій антені (задані $k_z$ і~поперечне $k_\perp\approx3{,}83/a$) густина плазми, в~якій
хвиля резонансно \enquote{вписується} у~трубку радіуса $a$, пропорційна полю $B$. Геліконні
джерела дають рекордні для ВЧ-розрядів густини $\sim10^{13}~\mathrm{см}^{-3}$ (задача~\ref{prb:Helicon}).

\begin{table}[ht]
	\centering\small
	\caption{Типові параметри технологічних джерел плазми (аргон).}
	\label{tab:rf-sources}
	\begin{platbl}{colspec={X[1.6,l,m]X[c,m]X[c,m]X[c,m]X[1.4,c,m]}}
		Джерело & $p$, мТорр & $n$, $\mathrm{см}^{-3}$ & $T_e$, еВ & Енергія іонів \\
		Тліючий (пост. струм) & $10^{2}$--$10^{4}$ & $10^{9}$--$10^{11}$ & $1$--$5$ & до $V_c$ \\
		ЄЗП ($13{,}56$~МГц) & $10$--$10^{3}$ & $10^{9}$--$10^{10}$ & $1$--$5$ & $10^{2}$--$10^{3}$~еВ, зв'язана з~$n$ \\
		ІЗП & $1$--$50$ & $10^{11}$--$10^{12}$ & $2$--$5$ & окреме зміщення \\
		ЕЦР ($2{,}45$~ГГц) & $0{,}1$--$10$ & $10^{11}$--$10^{12}$ & $2$--$7$ & окреме зміщення \\
		Геліконне & $1$--$10$ & $10^{12}$--$10^{13}$ & $3$--$6$ & окреме зміщення \\
	\end{platbl}
\end{table}

%% ========================================================
\section{Плазмова обробка матеріалів}\index{Плазмова обробка матеріалів|textbf}\label{sec:plasma-processing}
%% ========================================================

Мікросхему з~мільярдами транзисторів неможливо виготовити без плазми: близько третини
технологічних кроків є плазмовими. Плазма постачає на поверхню два \enquote{інструменти}:
потік \emph{хімічно активних радикалів} (атоми F, Cl, O, утворені дисоціацією молекул електронним
ударом) і~\emph{спрямований потік іонів} з~енергією, яку задає шар.

\emph{Травлення.} Атоми фтору травлять кремній і~самі по собі ($\mathrm{Si}+4\mathrm F\to\mathrm{SiF_4}\uparrow$),
але \emph{ізотропно}: маска підтравлюється збоку так само, як вглиб. Іони аргону з~енергією
$\sim500$~еВ розпорошують кремній, але повільно. Разом вони дають швидкість травлення, у~кілька
(до десяти) разів більшу за суму окремих швидкостей: іонне бомбардування руйнує поверхневий шар і~прискорює
утворення та десорбцію летких продуктів. Цей \emphx{іонно-стимульований}\index{Травлення!іонно-стимульоване|textbf} синергізм діє лише там, куди падають
іони, тобто на дно канавки, а~не на її стінки, бо пучок майже нормальний (кут $\sim1^\circ$, розд.~\ref{sec:discharge-sheath}).
Стінки часто ще й~захищає полімерна плівка з~фторвуглецевих радикалів ($\mathrm{CF_x}$), яку
іони здирають лише з~дна. Так виходять \emph{анізотропні} канавки з~відношенням глибини до ширини
понад $50{:}1$.

\emph{Осадження.} У~плазмохімічному осадженні з~газової фази (PECVD) енергію для дисоціації
прекурсора ($\mathrm{SiH_4}$, $\mathrm{NH_3}$) дають гарячі електрони, а~не нагрів підкладки. Тому
плівки аморфного кремнію чи нітриду кремнію ростуть при $\sim300\,^\circ\mathrm{C}$ замість $\gtrsim700\,^\circ\mathrm{C}$ у~
термічному процесі, і~це сумісно з~алюмінієвими контактами та скляними підкладками сонячних батарей
і~дисплеїв.

\emph{Розпорошення.} У~\emphx{магнетроні}\index{Магнетрон|textbf} над катодом-мішенню створюють арочне магнітне поле
$B\sim300$~Гс. Електрони ($\rho_e\sim0{,}1$~мм) замагнічені й~дрейфують у~схрещених полях~\eqref{eq:ExB}
по замкненій доріжці над мішенню, тому довго іонізують газ. Іони з~ларморівським радіусом
у~кілька сантиметрів практично не відчувають поля і~бомбардують мішень. Так можна підтримувати
густий розряд при тиску $\sim1$~мТорр, а~розпорошені атоми долітають до підкладки без зіткнень.

\emph{Розряди атмосферного тиску.} Діелектричний бар'єрний розряд (ДБР) між електродами,
вкритими діелектриком, складається з~багатьох наносекундних мікророзрядів: діелектрик
заряджається й~гасить кожен з~них раніше, ніж той перейде в~дугу. Газ лишається холодним, а~електрони
мають кілька еВ. На цьому ґрунтуються промислові озонатори, обробка полімерів, плазмова медицина.
Люмінесцентні та газорозрядні лампи~--- позитивний стовп, у~якому до $60~\%$ потужності
перетворюється на резонансне випромінювання (лінія Hg $254$~нм).

%% ========================================================
\section{Як це вимірюють: оптична емісія і~ВЧ-зонди}\label{sec:discharge-diagnostics}\index{Діагностика плазми!ВЧ-зонди}
%% ========================================================

\emph{Оптична емісійна спектроскопія} (ОЕС) найпростіша: розряд світиться сам, і~досить
спектрометра біля вікна камери. Інтенсивність лінії, збудженої прямим електронним ударом з~основного стану,
$I_X\propto n_en_Xk_X^{\rm exc}(T_e)$, але і~$n_e$, і~$k^{\rm exc}$ невідомі. Метод
\emphx{актинометрії}\index{Діагностика плазми!оптична емісія}\index{Актинометрія|textbf} обходить цю проблему: до газу додають кілька відсотків інертного газу
(\enquote{актинометра}, зазвичай Ar) відомої концентрації і~вибирають пару ліній із~близькими
порогами збудження, наприклад F~$703{,}7$~нм ($14{,}5$~еВ) і~Ar~$750{,}4$~нм ($13{,}5$~еВ). Тоді
відношення констант збудження майже не залежить від $T_e$, і
% ~~~~~~~~~~~~~~~~~~~~~~~~~~~~~~~~~~~~~~~~~~~~~~~~~~~~~~~~
\begin{equation}\label{eq:actinometry}
	\frac{n_{\mathrm F}}{n_{\mathrm{Ar}}}=C_{\mathrm{F,Ar}}\,\frac{I_{\mathrm F}}{I_{\mathrm{Ar}}},
\end{equation}
% ~~~~~~~~~~~~~~~~~~~~~~~~~~~~~~~~~~~~~~~~~~~~~~~~~~~~~~~~
де $C$~--- стала калібрування. На практиці зміну інтенсивності ліній продуктів реакції (наприклад,
CO або SiF) використовують для визначення моменту, коли травлення доходить до нижнього шару
(контроль закінчення процесу). Відношення інтенсивностей ліній з~\emph{різними} порогами,
навпаки, чутливе до хвоста розподілу і~дає оцінку $T_e$.

\emph{Зонд Ленгмюра у~ВЧ-плазмі.} Класична вольт-амперна характеристика зонда
(рис.~\ref{pic:probe-iv}) передбачає сталий потенціал плазми. У~ВЧ-розряді потенціал плазми
коливається з~амплітудою $\tilde V$ від кількох до сотень вольт. Через експоненту в~електронній гілці
усереднений струм зміщується і~розмивається: зонд показує завищену $T_e$ і~занижений плаваючий
потенціал. \emph{ВЧ-компенсований} зонд вирішує це двома способами: послідовно з~зондом вмикають
резонансні дроселі, налаштовані на $13{,}56$~МГц і~другу гармоніку, а~через ємність до зонда під'єднують
допоміжний електрод великої площі. Тоді зонд \enquote{плаває} разом з~потенціалом плазми, і~ВЧ-напруга на
його шарі стає малою, $e\tilde V_{\rm sh}\lesssim T_e/2$. З~компенсованої характеристики знаходять
не лише $T_e$ і~$n_e$, а~й~саму функцію розподілу за енергіями за \emph{методом Дрювестейна}:
% ~~~~~~~~~~~~~~~~~~~~~~~~~~~~~~~~~~~~~~~~~~~~~~~~~~~~~~~~
\begin{equation}\label{eq:druyvesteyn}
	g_e(\varepsilon)=\frac{2m}{e^2S}\sqrt{\frac{2e V}{m}}\,\frac{d^2I_e}{dV^2},\qquad \varepsilon=eV,
\end{equation}
% ~~~~~~~~~~~~~~~~~~~~~~~~~~~~~~~~~~~~~~~~~~~~~~~~~~~~~~~~
де $V$~--- гальмівний потенціал зонда відносно плазми, $S$~--- площа зонда, $g_e(\varepsilon)\,d\varepsilon$~--- число
електронів в~одиниці об'єму з~енергією в~інтервалі $d\varepsilon$. Формула випливає з~того, що струм на
зонд при гальмівному потенціалі $V$ несуть лише електрони з~енергією, більшою за $eV$, і~справедлива для
будь-якого ізотропного розподілу. Такі вимірювання показали, що в~ЄЗП низького тиску розподіл часто
\emph{двотемпературний}: холодні електрони в~центрі ($\sim0{,}5$~еВ) і~гарячі ($\sim3$~еВ), нагріті коливними
шарами. Для абсолютного значення $n_e$ незалежно від зонда використовують мікрохвильову
інтерферометрію або резонансний зонд-\enquote{шпильку}, власна частота якого зсувається на
$\omega_{pe}$.

\emph{Межі застосовності.} Актинометрія вимагає, щоб обидва верхні рівні збуджувалися прямим
електронним ударом з~основного стану (а~не дисоціативним збудженням молекул чи каскадом з~вищих рівнів)
і~розпадалися радіаційно, без гасіння зіткненнями. При тисках $\gtrsim1$~Торр і~в~молекулярних
газах ці умови порушуються, і~стала $C$ стає залежною від режиму. Зонд вносить у~плазму
збурення і~забруднюється плівками в~хімічно активних розрядах; друга похідна у~\eqref{eq:druyvesteyn}
підсилює шум, тож хвіст розподілу при $\varepsilon\gtrsim15$~еВ визначається з~великою похибкою. У~магнітному
полі (ЕЦР, магнетрон) струм електронів на зонд обмежено поперечною рухливістю, і~формула Дрювестейна
незастосовна.

\Needspace{8\baselineskip}
\begin{Summary}
\begin{itemize}
	\item Низькотемпературна плазма нерівноважна: $T_e\sim1$--$10$~еВ, $T_i\approx T_g$; її підтримує зовнішнє поле,
	а~іонізація компенсує втрати на стінки.
	\item Константа іонізації $k_{iz}\propto e^{-\varepsilon_{iz}/T_e}$~\eqref{eq:kiz} різко залежить від $T_e$,
	тому $T_e$ у~розрядах лежить у~вузькій смузі $2$--$5$~еВ.
	\item Пробій: лавина $\alpha=Ap\,e^{-Bp/E}$ плюс вторинна емісія, $\gamma(e^{\alpha d}-1)=1$; напруга пробою
	залежить лише від $pd$ (закон Пашена~\eqref{eq:paschen}) і~має мінімум $\sim200$--$300$~В.
	\item Тліючий розряд: майже вся напруга спадає в~катодному шарі; позитивний стовп підтримується
	балансом іонізації та амбіполярної дифузії (умова Шотткі~\eqref{eq:schottky}). Дуга має катодне
	падіння $\sim10$~В і~спадну характеристику.
	\item Потік іонів на поверхню $\approx0{,}6\,n_0u_B$ задає плазма; товщину високовольтного шару
	дає закон Чайлда--Ленгмюра~\eqref{eq:child-langmuir}, $s\sim r_D(eV_0/T_e)^{3/4}$.
	\item Глобальна модель: $T_e$ визначає баланс частинок, $k_{iz}/u_B=1/(n_gd_{\rm eff})$, незалежно від
	потужності; густина $n_0=P_{\rm abs}/(eu_BA_{\rm eff}\mathcal E_T)$ пропорційна потужності.
	\item ЄЗП зв'язує густину й~енергію іонів; ІЗП, ЕЦР і~геліконні джерела дають густішу плазму й~окреме
	керування енергією іонів через зміщення підкладки.
	\item Анізотропне травлення ґрунтується на синергії радикалів і~майже нормального потоку іонів із~шару.
	\item Діагностика: актинометрія для відносних концентрацій радикалів, ВЧ-компенсований зонд і~метод
	Дрювестейна для функції розподілу електронів.
\end{itemize}
\end{Summary}

\Needspace{9\baselineskip}
\section*{Задачі}
%% --------------------------------------------------------

\begin{problem}\label{prb:IonizationRate}\lvl{3}
	Нехай переріз іонізації поблизу порога $\sigma_{iz}=C(\varepsilon-\varepsilon_{iz})$ при $\varepsilon>\varepsilon_{iz}$.
	Усередніть $\sigma_{iz}v$ за максвеллівським розподілом і~отримайте~\eqref{eq:kiz}. У~скільки разів зросте
	$k_{iz}$ аргону ($\varepsilon_{iz}=15{,}8$~еВ) при зростанні $T_e$ від $2$ до $3$~еВ?
	\begin{solution}
		$k_{iz}=C\bar v_e(\varepsilon_{iz}+2T_e)e^{-\varepsilon_{iz}/T_e}$; зростання приблизно в~$\approx6$ разів
	(за апроксимацією з~розд.~\ref{sec:ionization-rate}~--- у~$\approx6{,}5$ раза).
	\end{solution}
\end{problem}

\begin{problem}\label{prb:Paschen}\lvl{2}
	Знайдіть положення і~величину мінімуму кривої Пашена~\eqref{eq:paschen}. Обчисліть їх для повітря
	($A=15~\mathrm{см}^{-1}\mathrm{Торр}^{-1}$, $B=365$~В/(см$\,\cdot\,$Торр), $\gamma=0{,}01$). Якою буде напруга
	пробою проміжку $d=1$~мм при атмосферному тиску?
	\begin{solution}
		$(pd)_{\min}=e\ln(1+1/\gamma)/A\approx0{,}84$~Торр$\,\cdot\,$см, $V_{\min}=B(pd)_{\min}\approx305$~В;
	при $pd=76$~Торр$\,\cdot\,$см $V_b\approx4{,}6$~кВ.
	\end{solution}
\end{problem}

\begin{problem}\label{prb:ChildLangmuir}\lvl{2}
	Виведіть~\eqref{eq:sheath-thickness} з~\eqref{eq:child-langmuir} і~умови $J=en_su_B$. Знайдіть товщину шару
	біля електрода під напругою $V_0=300$~В в~аргоновій плазмі з~$T_e=3$~еВ, $n_s=10^{10}~\mathrm{см}^{-3}$
	і~час прольоту іона через шар. Порівняйте його з~періодом поля $13{,}56$~МГц.
	\begin{solution}
		$r_D\approx0{,}13$~мм, $s\approx3{,}2$~мм; кінцева швидкість іона $\approx3{,}8\times10^6$~см/с,
	час прольоту $3s/v\approx0{,}25$~мкс, тобто $\approx3{,}4$ періоду поля.
	\end{solution}
\end{problem}

\begin{problem}\label{prb:GlobalModel}\lvl{3}
	Аргонова плазма при $p=10$~мТорр, $d_{\rm eff}=10$~см, $T_g=300$~К. Користуючись апроксимацією
	$k_{iz}\approx2{,}3\times10^{-8}T_e^{0{,}59}e^{-17{,}4/T_e}~\mathrm{см}^3/\mathrm{с}$, знайдіть $T_e$ з~\eqref{eq:global-te}.
	Яку густину підтримує поглинена потужність $500$~Вт при $A_{\rm eff}=10^3~\mathrm{см}^2$ і~$\mathcal E_T=80$~еВ?
	Як зміняться $T_e$ і~$n_0$, якщо потужність подвоїти, а~тиск зменшити вдвічі?
	\begin{solution}
		$T_e\approx2{,}8$~еВ, $n_0\approx1{,}5\times10^{11}~\mathrm{см}^{-3}$; від потужності
	$T_e$ не залежить, а~зменшення тиску вдвічі підвищує її до $\approx3{,}3$~еВ; густина зросте менш ніж
	удвічі, бо зростають $u_B$ і~$\mathcal E_T$ ($A_{\rm eff}$ вважаємо сталою).
	\end{solution}
\end{problem}

\begin{problem}\label{prb:ECRSource}\lvl{1}
	Знайдіть резонансне магнітне поле ЕЦР-джерела з~частотою $2{,}45$~ГГц і~критичну густину для звичайної
	хвилі цієї частоти. Чому мікрохвилі вводять з~боку сильного поля?
	\begin{solution}
		$B=2\pi fmc/e\approx875$~Гс; $n_c\approx7{,}4\times10^{10}~\mathrm{см}^{-3}$; при $\omega<\omega_{ce}$
	права хвиля є свистовою і~не має відсічки за густиною.
	\end{solution}
\end{problem}

\begin{problem}\label{prb:Helicon}\lvl{2}
	Геліконне джерело: трубка радіуса $a=2{,}5$~см, $f=13{,}56$~МГц, антена задає $\lambda_z=20$~см, $B=500$~Гс.
	Оцініть густину плазми з~\eqref{eq:helicon}, взявши $k_\perp=3{,}83/a$.
	\begin{solution}
		$k_\perp\approx1{,}5~\mathrm{см}^{-1}$, $k_z\approx0{,}31~\mathrm{см}^{-1}$, $k\approx1{,}6~\mathrm{см}^{-1}$;
	$n\approx1{,}4\times10^{13}~\mathrm{см}^{-3}$.
	\end{solution}
\end{problem}

\begin{problem}\label{prb:Magnetron}\lvl{2}
	У~магнетроні $B=300$~Гс, енергія електронів $\sim3$~еВ, іонів аргону~--- до $\sim10$~еВ. Порівняйте
	ларморівські радіуси електронів та іонів із~товщиною розрядної області $\sim1$~см.
	\begin{solution}
		$\rho_e\approx0{,}15$~мм (замагнічені), $\rho_i\approx7$~см (незамагнічені).
	\end{solution}
\end{problem}

%% --------------------------------------------------------

\Needspace{6\baselineskip}
\begin{Questions}
\begin{quizlist}
	\item Чому формула Саха незастосовна до плазми тліючого розряду, але застосовна до дуги атмосферного тиску?
	\item Поясніть фізично, чому крива Пашена має мінімум.
	\item Чому в~глобальній моделі температура електронів не залежить від потужності, а~густина пропорційна їй?
	\item Чому потік іонів на поверхню не залежить від прикладеної до неї напруги, а~товщина шару залежить?
	\item Звідки береться автозміщення в~ємнісному розряді і~чому пластину кладуть на менший електрод?
	\item У~чому перевага індуктивного джерела над ємнісним для травлення?
	\item Чому для актинометрії вибирають лінії з~близькими порогами збудження?
\end{quizlist}
\end{Questions}

\Closesolutionfile{answer}
